\documentclass[11pt,a4paper]{article}
\usepackage[utf8]{inputenc}
\usepackage[french]{babel}
\usepackage[T1]{fontenc}
\usepackage{geometry}
\geometry{top=2.5cm, bottom=2.5cm, left=2.2cm, right=2.2cm}
\usepackage{xcolor}
\usepackage{hyperref}
\usepackage{listings}
\usepackage{tcolorbox}
\usepackage{tikz}
\usetikzlibrary{shapes.geometric, arrows.meta, positioning, shadows}
\usepackage{array}
\usepackage{booktabs}
\usepackage{fancyhdr}
\usepackage{float}

% Palette de couleurs professionnelle
\definecolor{navyprimary}{RGB}{15, 23, 42}
\definecolor{blueaccent}{RGB}{29, 78, 216}
\definecolor{slategray}{RGB}{71, 85, 105}
\definecolor{codebg}{RGB}{248, 250, 252}
\definecolor{codeframe}{RGB}{226, 232, 240}
\definecolor{greenaccent}{RGB}{22, 101, 52}
\definecolor{purplaccent}{RGB}{126, 34, 206}

% Configuration de hyperref
\hypersetup{
    colorlinks=true,
    linkcolor=blueaccent,
    filecolor=purplaccent,
    urlcolor=blueaccent,
    pdftitle={Rapport d'Architecture IncidentFlow},
    pdfpagemode=FullScreen
}

% Configuration de headheight pour fancyhdr
\setlength{\headheight}{14.5pt}
\pagestyle{fancy}
\fancyhf{}
\rhead{\textcolor{slategray}{\small IncidentFlow --- Architecture \& Flux de Données}}
\lhead{\textcolor{slategray}{\small Rapport Technique}}
\rfoot{\textcolor{slategray}{\small Page \thepage}}
\lfoot{\textcolor{slategray}{\small NetMar Technologies}}
\renewcommand{\headrulewidth}{0.4pt}

% Définition de JavaScript pour listings
\lstdefinelanguage{JavaScript}{
  keywords={break, case, catch, continue, debugger, default, delete, do, else, finally, for, function, if, in, instanceof, new, return, switch, this, throw, try, typeof, var, void, while, with, let, const, async, await, import, export, from},
  morecomment=[l]{//},
  morecomment=[s]{/*}{*/},
  morestring=[b]',
  morestring=[b]"
}

% Configuration de listings pour le code
\lstset{
    backgroundcolor=\color{codebg},
    frame=single,
    rulecolor=\color{codeframe},
    basicstyle=\ttfamily\footnotesize\color{navyprimary},
    keywordstyle=\bfseries\color{blueaccent},
    commentstyle=\itshape\color{slategray},
    stringstyle=\color{greenaccent},
    breaklines=true,
    captionpos=b,
    showstringspaces=false,
    tabsize=2
}

\begin{document}

% Page de titre
\begin{titlepage}
    \centering
    \vspace*{1.5cm}
    
    {\Huge \bfseries \textcolor{navyprimary}{IncidentFlow}}\\[0.4cm]
    {\Large \textcolor{blueaccent}{Spécification Détaillée de l'Architecture \& du Flux de Données}}\\[0.2cm]
    {\large \textcolor{slategray}{Des Entités Base de Données à l'Interface Utilisateur Real-Time}}\\[1.5cm]
    
    \begin{tcolorbox}[colback=blueaccent!5,colframe=blueaccent,arc=6pt,title=\textbf{Synthèse du Document}]
    Ce document présente le fonctionnement interne du système \textbf{IncidentFlow}. Il détaille le parcours complet des données relatif aux \textbf{Utilisateurs}, aux \textbf{Rôles}, et au \textbf{Workflow Dynamique} (États et Transitions), depuis la persistance en base de données PostgreSQL jusqu'à l'affichage et la synchronisation en temps réel sur le frontend React.
    \end{tcolorbox}

    \vfill
    
    \begin{tabular}{ll}
        \textbf{Application :} & IncidentFlow (Gestion d'Incidents Interne) \\
        \textbf{Stack Backend :} & Java 17, Spring Boot 3, JPA / Hibernate, PostgreSQL \\
        \textbf{Stack Frontend :} & React 19, Vite, TailwindCSS / Vanilla CSS \\
        \textbf{Auteur :} & Équipe de Développement NetMar \\
        \textbf{Date :} & \today \\
    \end{tabular}
    
    \vspace*{0.8cm}
\end{titlepage}

\tableofcontents
\newpage

% ---------------------------------------------------------
% SECTION 1: VUE D'ENSEMBLE DE L'ARCHITECTURE
% ---------------------------------------------------------
\section{Vue d'Ensemble de l'Architecture 3-Tiers}

L'application \textbf{IncidentFlow} repose sur une architecture robuste à trois tiers. Les composants interagissent de manière découplée via des API REST normées et du JSON.

\begin{figure}[H]
\centering
\begin{tikzpicture}[node distance=1.8cm, auto, >=Stealth, 
    layer/.style={rectangle, draw=navyprimary, fill=codebg, rounded corners=6pt, minimum width=14cm, minimum height=2.2cm, thick},
    component/.style={rectangle, draw=blueaccent, fill=white, rounded corners=4pt, minimum width=3.6cm, minimum height=1.2cm, font=\small\bfseries, align=center}]

    % Couche Frontend
    \node[layer, label=above:{\textbf{\textcolor{navyprimary}{Couche 1 : Frontend (React 19 / Vite)}}}] (frontlayer) {};
    \node[component] (reactui) at (frontlayer.west) [xshift=2.5cm] {Composants UI\\(\texttt{App.jsx})};
    \node[component] (reactstate) at (frontlayer.center) {Gestion d'État\\(\texttt{useState}, \texttt{useEffect})};
    \node[component] (fetchapi) at (frontlayer.east) [xshift=-2.5cm] {Client API Fetch\\(\texttt{getHeaders()})};

    % Couche Backend
    \node[layer, label=above:{\textbf{\textcolor{navyprimary}{Couche 2 : Backend (Spring Boot 3)}}}, below=1.5cm of frontlayer] (backlayer) {};
    \node[component] (controller) at (backlayer.west) [xshift=2.5cm] {Contrôleurs REST\\(\texttt{WorkflowController})};
    \node[component] (service) at (backlayer.center) {Couche Métier\\(\texttt{WorkflowService})};
    \node[component] (repo) at (backlayer.east) [xshift=-2.5cm] {JPA Repositories\\(\texttt{WorkflowRepository})};

    % Couche BDD
    \node[layer, label=above:{\textbf{\textcolor{navyprimary}{Couche 3 : Base de Données (PostgreSQL)}}}, below=1.5cm of backlayer] (dblayer) {};
    \node[component] (dbusers) at (dblayer.west) [xshift=2.5cm] {Table \texttt{users}\\+ \texttt{roles}};
    \node[component] (dbwf) at (dblayer.center) {Table \texttt{workflows}};
    \node[component] (dbstates) at (dblayer.east) [xshift=-2.5cm] {Tables \texttt{workflow\_states}\\+ \texttt{transitions}};

    % Flèches de communication
    \draw[<->, very thick, color=blueaccent] (fetchapi) -- node[right, font=\footnotesize] {HTTP REST / JSON} (controller);
    \draw[<->, very thick, color=blueaccent] (controller) -- (service);
    \draw[<->, very thick, color=blueaccent] (service) -- (repo);
    \draw[<->, very thick, color=blueaccent] (repo) -- node[right, font=\footnotesize] {JPA / SQL / Hibernate} (dblayer);

\end{tikzpicture}
\caption{Diagramme de flux de données de bout en bout de l'application IncidentFlow}
\end{figure}

% ---------------------------------------------------------
% SECTION 2: COUCHE BASE DE DONNÉES ET MODÈLE OBJET JPA
% ---------------------------------------------------------
\section{Couche 1 : Base de Données \& Modèle Objet JPA}

La persistance des données s'appuie sur le SGBD \textbf{PostgreSQL}. La modélisation est représentée côté Java par des entités JPA (Java Persistence API) annotées avec Hibernate.

\subsection{Gestion des Utilisateurs et des Rôles}
Les utilisateurs et les rôles sont modélisés par deux entités liées par une relation de type \texttt{@ManyToOne}.

\begin{itemize}
    \item \textbf{Fichier Rôle :} \texttt{backend/src/main/java/com/netmar/incidentflow/model/Role.java}
    \item \textbf{Fichier Utilisateur :} \texttt{backend/src/main/java/com/netmar/incidentflow/model/User.java}
\end{itemize}

\begin{lstlisting}[language=Java, caption=Structure JPA de l'entité User et sa relation avec Role]
@Entity
@Table(name = "users")
public class User {
    @Id @GeneratedValue(strategy = GenerationType.IDENTITY)
    private Long id;
    
    @Column(nullable = false, unique = true)
    private String email;
    
    private String name;
    
    @ManyToOne
    @JoinColumn(name = "role_id")
    private Role role; // Contient "Administrateur", "Responsable", "Operateur", etc.
}
\end{lstlisting}

\subsection{Modélisation du Workflow Dynamique (États et Transitions)}
Le moteur de workflow est composé de trois entités en cascade : \texttt{Workflow}, \texttt{WorkflowState} et \texttt{WorkflowTransition}.

\begin{itemize}
    \item \textbf{Fichier d'Entité :} \texttt{backend/src/main/java/com/netmar/incidentflow/model/Workflow.java}
    \item \textbf{Fichier États :} \texttt{backend/src/main/java/com/netmar/incidentflow/model/WorkflowState.java}
    \item \textbf{Fichier Transitions :} \texttt{backend/src/main/java/com/netmar/incidentflow/model/WorkflowTransition.java}
\end{itemize}

\begin{lstlisting}[language=Java, caption=Association en Cascade et Orphelins du Workflow]
@Entity
@Table(name = "workflows")
public class Workflow {
    @Id @GeneratedValue(strategy = GenerationType.IDENTITY)
    private Long id;
    private String name;
    
    @OneToMany(mappedBy = "workflow", cascade = CascadeType.ALL, orphanRemoval = true, fetch = FetchType.EAGER)
    private List<WorkflowState> states = new ArrayList<>();

    @OneToMany(mappedBy = "workflow", cascade = CascadeType.ALL, orphanRemoval = true, fetch = FetchType.EAGER)
    private List<WorkflowTransition> transitions = new ArrayList<>();
}
\end{lstlisting}

\subsection{Initialisation des Données en Base (\texttt{DataInitializer.java})}
Au démarrage du serveur backend, la classe d'initialisation vérifie l'existence des rôles et du workflow standard.

\begin{itemize}
    \item \textbf{Fichier Source :} \texttt{backend/src/main/java/com/netmar/incidentflow/config/DataInitializer.java}
\end{itemize}

Si aucun workflow n'est présent (\texttt{workflowRepository.count() == 0}), un workflow par défaut nommé \textbf{Workflow Standard} est créé avec 5 états principaux (\textit{Nouveau}, \textit{Assigné}, \textit{En cours}, \textit{Résolu}, \textit{Clôturé}) et leurs transitions respectives.

% ---------------------------------------------------------
% SECTION 3: COUCHE BACKEND (SPRING BOOT 3)
% ---------------------------------------------------------
\section{Couche 2 : Couche Métier \& REST Backend (Spring Boot 3)}

La couche backend assure la sécurité, la validation métier, le traitement des entités détachées JPA et l'exposition des API REST.

\subsection{Accès aux Données (Repositories)}
Les interfaces Spring Data JPA permettent d'exécuter les requêtes SQL nécessaires :
\begin{itemize}
    \item \texttt{UserRepository.java} : \texttt{findByEmail(String email)}
    \item \texttt{RoleRepository.java} : \texttt{findByName(String name)}
    \item \texttt{WorkflowRepository.java} : \texttt{findByCategory(String category)}
\end{itemize}

\subsection{Couche Services \& Gestion des Entités Détachées JPA}
La logique métier est centralisée dans \texttt{UserService.java} et \texttt{WorkflowService.java}.

\subsubsection{Résolution du Bug Détaché JPA dans \texttt{WorkflowService.java}}
Lorsqu'un utilisateur modifie un workflow depuis le frontend, les objets JSON reçus contiennent des identifiants existants d'états mais ne sont pas rattachés à la session Hibernate courante. Un \texttt{repository.save()} direct provoquait l'erreur \texttt{detached entity passed to persist}.

\begin{itemize}
    \item \textbf{Fichier Modifié :} \texttt{backend/src/main/java/com/netmar/incidentflow/service/WorkflowService.java}
\end{itemize}

\begin{lstlisting}[language=Java, caption=Résolution de la persistance des entités détachées]
@Transactional
@CacheEvict(value = {"workflows", "workflows-category"}, allEntries = true)
public Workflow saveWorkflow(Workflow workflow) {
    if (workflow.getId() != null) {
        // 1. Charger l'entite geree par JPA
        Workflow existing = workflowRepository.findById(workflow.getId()).orElseThrow();
        existing.setName(workflow.getName());
        
        // 2. Vider les collections (declenche orphanRemoval = true)
        existing.getStates().clear();
        existing.getTransitions().clear();
        
        // 3. Re-instancier et rattacher les nouveaux enfants a l'entite geree
        if (workflow.getStates() != null) {
            for (WorkflowState s : workflow.getStates()) {
                existing.getStates().add(WorkflowState.builder()
                        .name(s.getName()).label(s.getLabel())
                        .colorClass(s.getColorClass()).workflow(existing).build());
            }
        }
        return workflowRepository.save(existing);
    }
    return workflowRepository.save(workflow);
}
\end{lstlisting}

\subsection{Contrôleurs REST (API Endpoints)}
Les données sont exposées au frontend via les contrôleurs REST :

\begin{table}[H]
\centering
\small
\begin{tabular}{llll}
\toprule
\textbf{Méthode} & \textbf{Endpoint API} & \textbf{Contrôleur Backend} & \textbf{Description} \\
\midrule
\texttt{GET} & \texttt{/api/users} & \texttt{UserController.java} & Récupère tous les utilisateurs \\
\texttt{GET} & \texttt{/api/roles} & \texttt{UserController.java} & Récupère la liste des rôles \\
\texttt{GET} & \texttt{/api/workflows} & \texttt{WorkflowController.java} & Récupère le workflow unique \\
\texttt{PUT} & \texttt{/api/workflows/\{id\}} & \texttt{WorkflowController.java} & Enregistre les modifications du workflow \\
\texttt{DELETE} & \texttt{/api/incidents/\{code\}} & \texttt{IncidentController.java} & Supprime un incident (Admin requis) \\
\bottomrule
\end{tabular}
\caption{Principaux points de terminaison REST de l'application}
\end{table}

% ---------------------------------------------------------
% SECTION 4: COUCHE FRONTEND (REACT 19 / VITE)
% ---------------------------------------------------------
\section{Couche 3 : Interface Utilisateur Frontend (React 19 / Vite)}

Le frontend React consomme l'API REST via la fonction système \texttt{fetch} et met à jour dynamiquement l'interface grâce au système d'états réactifs de React.

\subsection{Chargement Initial et Reconstitution des États (\texttt{App.jsx})}
Au chargement du composant principal \texttt{App.jsx}, le hook \texttt{useEffect} déclenche les requêtes asynchrones pour récupérer les workflows, les utilisateurs et les incidents.

\begin{itemize}
    \item \textbf{Fichier Frontend Principal :} \texttt{frontend/src/App.jsx}
\end{itemize}

\begin{lstlisting}[language=JavaScript, caption=Appels asynchrones et synchronisation d'état React]
const [workflows, setWorkflows] = useState([]);
const [activeWorkflow, setActiveWorkflow] = useState(null);

const fetchWorkflows = async () => {
  try {
    const res = await fetch(`${API_BASE}/workflows`, { headers: getHeaders() });
    if (res.ok) {
      const data = await res.json();
      setWorkflows(data);
      if (data.length > 0) setActiveWorkflow(data[0]);
    }
  } catch (err) { console.error(err); }
};
\end{lstlisting}

\subsection{Propagation en Temps Réel dans l'Interface (Real-Time UI Sync)}
Chaque modification du workflow (ajout/suppression d'un état ou d'une transition) est immédiatement propagée sur trois composants clés de l'interface :

\begin{enumerate}
    \item \textbf{Éditeur Graphique de Workflow :} L'administrateur modifie les nœuds et les connexions dans l'éditeur interactif.
    \item \textbf{Filtre de Statut Dynamique dans la Liste des Incidents :} Le menu déroulant de sélection du statut s'alimente directement depuis \texttt{activeWorkflow.states} :
\begin{lstlisting}[language=JavaScript]
<select value={statusFilter} onChange={(e) => setStatusFilter(e.target.value)}>
  <option value="Tous">Tous</option>
  {activeWorkflow?.states.map(s => (
    <option key={s.id || s.name} value={s.name}>{s.label || s.name}</option>
  ))}
</select>
\end{lstlisting}
    \item \textbf{Stepper Horizontal \& Boutons d'Action (Fiche de Détail) :} Le stepper visuel d'avancement ainsi que les boutons d'action autorisés (\textit{Passer à : En cours}, etc.) s'évaluent dynamiquement en fonction des états du workflow actif :
\begin{lstlisting}[language=JavaScript]
const currentWf = activeWorkflow || workflows[0];
// Rendu dynamique du Stepper Horizontal et des Actions Disponibles
\end{lstlisting}
\end{enumerate}

% ---------------------------------------------------------
% SECTION 5: SCÉNARIO COMPLET ÉTAPE PAR ÉTAPE
% ---------------------------------------------------------
\section{Scénario Complet : De la Modification en Base au Rendu Temps Réel}

Pour illustrer le passage complet des données à votre encadrant, voici les 6 étapes de l'exécution complète lorsqu'un Administrateur ajoute un état au workflow :

\begin{tcolorbox}[colback=greenaccent!5,colframe=greenaccent,arc=4pt,title=\textbf{Parcours d'une Modification en 6 Étapes}]
\begin{enumerate}
    \item \textbf{Action Utilisateur (Frontend) :} L'administrateur ajoute un nouvel état (ex: \textit{En Validation}) dans la page de gestion du Workflow et clique sur \textit{"Enregistrer les modifications"}.
    \item \textbf{Requête HTTP PUT (Client REST) :} React envoie une requête \texttt{PUT /api/workflows/1} avec le payload JSON structuré contenant les nouveaux états et transitions.
    \item \textbf{Réception \& Sécurité (Backend) :} \texttt{WorkflowController.java} reçoit la requête, transmet l'objet à \texttt{WorkflowService.java} et vérifie le rôle de l'utilisateur.
    \item \textbf{Traitement JPA \& Persistance SQL (Service \& Database) :}
    \begin{itemize}
        \item \texttt{WorkflowService} charge l'entité gérée \texttt{existing}.
        \item Effacement des anciennes collections via \texttt{.clear()} (déclenchant les \texttt{DELETE SQL}).
        \item Insertion des nouveaux états via \texttt{save()} (déclenchant les \texttt{INSERT SQL} dans \texttt{workflow\_states}).
    \end{itemize}
    \item \textbf{Réponse HTTP 200 \& Mise à jour du State React :} Le backend renvoie l'objet \texttt{Workflow} mis à jour. \texttt{App.jsx} met à jour son état React local via \texttt{setActiveWorkflow(updated)} et déclenche un \texttt{fetchIncidents()}.
    \item \textbf{Mise à Jour Visuelle Temps Réel (DOM React) :} React ré-évalue le composant. Le nouvel état \textit{"En Validation"} apparaît instantanément dans :
    \begin{itemize}
        \item Le menu déroulant du filtre de statut de la page d'incidents.
        \item Le stepper horizontal de progression de la fiche détaillée d'incident.
        \item Le panneau des actions de transition disponibles.
    \end{itemize}
\end{enumerate}
\end{tcolorbox}

% ---------------------------------------------------------
% SECTION 6: TABLEAU RÉCAPITULATIF DES FICHIERS
% ---------------------------------------------------------
\section{Tableau Récapitulatif des Fichiers du Projet}

\begin{table}[H]
\centering
\small
\begin{tabular}{lll}
\toprule
\textbf{Composant / Fichier} & \textbf{Emplacement dans le Projet} & \textbf{Rôle Principal dans le Flux} \\
\midrule
\texttt{User.java} & \texttt{backend/.../model/User.java} & Entité JPA décrivant l'utilisateur \\
\texttt{Role.java} & \texttt{backend/.../model/Role.java} & Entité JPA décrivant les rôles/permissions \\
\texttt{Workflow.java} & \texttt{backend/.../model/Workflow.java} & Entité racine du moteur de workflow \\
\texttt{WorkflowState.java} & \texttt{backend/.../model/WorkflowState.java} & Entité enfant décrivant chaque état \\
\texttt{WorkflowTransition.java} & \texttt{backend/.../model/WorkflowTransition.java} & Entité enfant décrivant les transitions \\
\texttt{DataInitializer.java} & \texttt{backend/.../config/DataInitializer.java} & Script de seeding initial au démarrage \\
\texttt{UserService.java} & \texttt{backend/.../service/UserService.java} & Gestion métier utilisateurs/rôles/auth \\
\texttt{WorkflowService.java} & \texttt{backend/.../service/WorkflowService.java} & Logique de persistance et re-attachement \\
\texttt{WorkflowController.java} & \texttt{backend/.../controller/WorkflowController.java} & Endpoints REST pour le workflow \\
\texttt{IncidentController.java} & \texttt{backend/.../controller/IncidentController.java} & Endpoints REST incidents + suppression \\
\texttt{App.jsx} & \texttt{frontend/src/App.jsx} & Interface React, état global et rendu UI \\
\bottomrule
\end{tabular}
\caption{Matrice des fichiers clés du projet IncidentFlow}
\end{table}

\end{document}
